\subsection{Kvantiniai vartai}
\label{sec:kvantiniai-vartai}

Kvantiniai algoritmai sudaromi iš vartų, kurie keičia kvantinio registro amplitudes. Teoriškai kvantiniai vartai aprašomi unitariomis matricomis, tačiau šiame darbe svarbiausia jų programinė semantika: kiekvienas vartas yra operacija, kurią galima pritaikyti vienam ar keliems kubitams.

Groverio algoritmo realizacijoje naudojami šie vartai:

\begin{itemize}
    \item \textbf{Pauli-X vartai} -- veikia kaip kvantinis NOT: $\lvert 0 \rangle$ pakeičia į $\lvert 1 \rangle$, o $\lvert 1 \rangle$ pakeičia į $\lvert 0 \rangle$. Programoje jie naudojami pažymėtos būsenos paruošimui oracle dalyje.
    \item \textbf{Hadamard vartai} -- sukuria superpoziciją. Pritaikius šį vartą $\lvert 0 \rangle$ būsenai, gaunama vienoda $\lvert 0 \rangle$ ir $\lvert 1 \rangle$ amplitudžių kombinacija:
    \begin{equation}
        H\lvert 0 \rangle = \frac{\lvert 0 \rangle + \lvert 1 \rangle}{\sqrt{2}}.
    \end{equation}
    Pritaikius Hadamard vartus visiems paieškos registro kubitams, gaunama vienoda visų $2^n$ būsenų superpozicija.
    \item \textbf{CNOT vartai} -- dviejų kubitų vartai, kuriuose vienas kubitas yra kontrolinis, o kitas -- taikinys. Taikinys apverčiamas tik tada, kai kontrolinis kubitas yra būsenoje $1$.
    \item \textbf{Toffoli vartai} -- trijų kubitų vartai, dar vadinami kontroliuojamu kontroliuojamu NOT. Taikinys apverčiamas tik tada, kai abu kontroliniai kubitai yra $1$. Šis vartas naudojamas kelių sąlygų oracle ir difuzijos operatoriaus konstrukcijai.
\end{itemize}
